\begin{abstract}
Pedestrian Attribute Recognition (PAR) in surveillance applications is severely challenged by cross-camera domain shifts, extreme class imbalance between common and rare attributes, and spatial misalignment caused by pedestrian pose and viewpoint variations. In this work, we propose an integrated framework developed for the UPAR Challenge 2027 (Track 1: Pedestrian Attribute Recognition) consisting of \textbf{Anatomical Strip-Routed ConvNeXt with Source-Prior Harmonic Calibration (ASR-HC)} and a \textbf{Joint Execution Vectorization (JEV)} shared-stem inference engine. Our approach builds upon a modern ConvNeXt backbone with frozen early stages to retain robust multi-domain representations while spatially partitioning deep feature maps into three anatomical horizontal strips: Head/Shoulders, Torso/Upper-Body, and Lower-Body/Footwear. An 8-head Cross-Strip Multi-Head Self-Attention module refines token interactions across body parts with learnable positional embeddings before routing localized features to their corresponding attribute classifiers. To counteract the long-tailed distribution and cross-color hallucination inherent to surveillance imagery, we formulate an Asymmetric Source-Prior Harmonic Calibrator with continuous odds-ratio probability warping and intra-group competitive suppression. Evaluated on the comprehensive 33,407 validation images of the harmonized UPAR benchmark (PA-100K, PETA, and Market1501), ASR-HC substantially elevates Instance Precision to \textbf{84.29\%--86.28\%} (+6.56\% to +8.55\% over baseline) and Instance Accuracy to \textbf{74.99\%}. Furthermore, our vectorized JEV dual-head ensemble achieves a state-of-the-art Harmonic Mean of \textbf{81.86\%} (80.59\% label mA and 83.18\% instance F1) with minimal computational overhead.
\end{abstract}
